\documentclass[fleqn]{scrartcl}
% \usepackage{amsfonts}
% mathtools loads
% - amsmath -> amsbsy, amsopn, amstext
% - kezval, calc, mhsetpu, graphicx
\usepackage{mathtools} % -> amsmath -> amsbsy, amsopn, amstext
\usepackage{tikz}
\usepackage[ngerman]{babel}
\usepackage[utf8]{inputenc}

\usetikzlibrary {arrows.meta, angles, quotes}
\usetikzlibrary{calc}
\usetikzlibrary{patterns}
\usetikzlibrary{positioning}

\tikzset{
    % pics/lager/.style n args = {1}{code={
    pics/lager/.style = {code={
    \coordinate (A)  at ( 90:2.0mm);
    \coordinate (B)  at (210:2.0mm);
    \coordinate (C)  at (-30:2.0mm);
    \coordinate (GL) at ([shift={(-1mm, #1)}] B.west);
    \coordinate (GR) at ([shift={( 1mm, #1)}] C.east);
    \coordinate (DL) at ([shift={( 0mm,-2mm)}] GL.south);
    \coordinate (DR) at ([shift={( 0mm,-2mm)}] GR.south);

    \draw (A) -- (B) -- (C) -- cycle;
    \draw [fill=white] (A) circle (0.5mm);
    \draw (GL) -- (GR);
    \pattern[pattern=north east lines] (GL) -- (GR) -- (DR) -- (DL) -- cycle;
    }}
}


% Numerierung der Formel innerhalb einer Sektion
\numberwithin{equation}{section}

% Abkürzung für mathbf
\def\*#1{\mathbf{#1}}

\title{Superposition \\
für Stefan Mattes}
\author{Janos Sarközi}

\begin{document}
\maketitle
\newpage
\section{Vorbereitungen}
Zuerst werden die Biegelinien verschiedener Konstruktionen betrachtet, die eine wesentliche Rolle bei der Ermittlung der Momente und Kräfte in statisch unbestimmten Systemen spielen werden. Dafür werden aus den Gleichungen
\[
    Q'=-q,\hspace{1ex}M'=Q,\hspace{1ex},\psi'=\frac{M}{EI},\hspace{1ex}w'=-\psi
\]
die Differentialgleichung der Biegeline in zwei Formen aufgeschrieben. Die erste Form ist
\[
    w''=-\frac{M}{EI}
\]
Dies fogt aus den letzten beiden Gleichungen. Setzt man die zweite Gleichung in die erte ein, so erhält man $M''=-q$ und damit
\[
    (EIw'')''=q
\]
Wird nun $EI$ als konstant angenommen, dann erhalten wir die finale Version der Biegelinie
:
\[
    EIw^{IV}=q
\]
\begin{figure}[h]
\centering
\begin{tikzpicture}
    \pic (LA) at (0.0, 0.0) {lager={ 0.0mm}};
    \pic (LB) at (4.0, 0.0) {lager={-0.7mm}};
    \node [draw, minimum width=4.4cm, above=0.0 and 0.0 of $(LAA)!0.5!(LBA)$] (z) {};
\end{tikzpicture}
\caption{Brücke}
\label{fig:paralel}
\end{figure}
\end{document}
